%
% ---
% id : "electrotechnique_knx_20260428_612"
% title : "Affirmations sur la technologie KNX"
% domain : "Électrotechnique"
% subdomain : "KNX"
% tags : ["knx", "bus", "topologie", "adresse physique", "tension bus", "ts"]
% difficulty : 1
% ---

\documentclass[a4paper,11pt,answers,addpoints]{exam}

% ---------------------------------------------------------
% LANGUE, ENCODAGE, FONTS
% ---------------------------------------------------------
\usepackage[utf8]{inputenc}

% ---------------------------------------------------------
% GÉOMÉTRIE ET MISE EN PAGE
% ---------------------------------------------------------
\usepackage[left=1.7cm,right=1.5cm,top=2cm,bottom=1cm]{geometry}
%\usepackage{a4wide}
\usepackage{microtype}      % Meilleure répartition du texte
\usepackage{float}          % Positionnement [H]
\usepackage{needspace}      % Saut conditionnel intelligent

% Éviter les veuves/orphelines (optionnel, à décommenter si besoin)
% \widowpenalty=10000
% \clubpenalty=10000

% Espacement vertical flexible
\raggedbottom
\setlength{\parskip}{0.5em plus 0.2em minus 0.1em}
\setlength{\parindent}{0pt}


% ---------------------------------------------------------
% MATHÉMATIQUES
% ---------------------------------------------------------
\usepackage{amsmath, amssymb, bm, esvect}

% Vecteurs
\newcommand{\V}{\overrightarrow}
\def\vect#1{\overrightarrow{\strut#1}}
\providecommand{\Degrees}[0]{\ensuremath{^\circ}}

% ---------------------------------------------------------
% UNITÉS ET GRANDEURS (SIunitx)
% ---------------------------------------------------------
\usepackage{siunitx}
\sisetup{output-decimal-marker={,}}
\DeclareSIUnit \var {var}
\DeclareSIUnit \VA {VA}
\DeclareSIUnit \kVA {kVA}
\DeclareSIUnit[quantity-product={}] \degree{\text{\textdegree}}

% ---------------------------------------------------------
% GRAPHIQUES : TikZ + CircuitikZ + PGFPlots
% ---------------------------------------------------------
\usepackage{tikz}
\usetikzlibrary{
	arrows.meta,
	intersections,
	decorations.pathreplacing,
	calligraphy,
	positioning
}

\usepackage[european,straightvoltages,EFvoltages]{circuitikz}

\usepackage{tkz-fct}
\usepackage{pgfplots}
\pgfplotsset{compat=1.12}

% Symbole étoile-triangle personnalisé
\newcommand{\wye}{%
	\mathbin{\tikz[x=1ex,y=1ex]{%
			\draw[line width=.1ex] (0,0)--(45:1)--++(-45:1) (45:1)--++(0,1);
	}}%
}


\usepackage{sankey}

\pgfdeclarelayer{background}
\pgfdeclarelayer{foreground}
\pgfdeclarelayer{sankeydebug}
\pgfsetlayers{background,main,foreground,sankeydebug}


% ---------------------------------------------------------
% TABLEAUX
% ---------------------------------------------------------
\usepackage{tabularx}
\usepackage{tabu}

% ---------------------------------------------------------
% HYPERLIENS
% ---------------------------------------------------------
\usepackage[colorlinks=true,
menucolor=red,
breaklinks=true,
urlcolor=blue,
linkcolor=blue]{hyperref}

% ---------------------------------------------------------
% COMMANDES PERSONNELLES
% ---------------------------------------------------------
\newcommand\nomauteur{bdminasmorgul@protonmail.com}
\newcommand\dateTe{26.04.2026}
\newcommand\brancheTe{TS}

% ---------------------------------------------------------
% STYLE DE L'EXERCICE
% ---------------------------------------------------------
\pagestyle{headandfoot}
\firstpageheadrule
\runningheadrule

\firstpageheader{\brancheTe\ (\nomauteur)}{Élève : }{(\dateTe) \thepage/\numpages}
\runningheader{\brancheTe\ (\nomauteur)}{Élève : }{(\dateTe) \thepage/\numpages}

% Compteur en bas de page
\newcommand{\totalpointtts}{%
	\begin{tikzpicture}[remember picture, overlay]
		\node[anchor=south west, inner sep=0pt, outer sep=0pt]
		at ([xshift=0.5cm,yshift=0.5cm]current page.south west)
		{\rotatebox{90}{Total des points : \numpoints}};
	\end{tikzpicture}
}

\firstpagefooter{\totalpointtts}{}{}
\runningfooter{\totalpointtts}{}{}

% Réglages exam
\printanswers
\addpoints
\pointsinmargin
\bracketedpoints

% ---------------------------------------------------------
% LISTINGS (code)
% ---------------------------------------------------------
\usepackage{xcolor}
\usepackage{listings}

\lstdefinestyle{octavestyle}{
	language=Matlab,
	basicstyle=\ttfamily\small,
	keywordstyle=\color{blue},
	commentstyle=\color{green!50!black},
	stringstyle=\color{red},
	stepnumber=1,
	frame=single,
	breaklines=true,
	breakatwhitespace=true,
	tabsize=4
}

% ---------------------------------------------------------
% DOCUMENT
% ---------------------------------------------------------
\begin{document}
	
	
	\unframedsolutions
	\pointsinmargin
	\bracketedpoints
	\section*{Préparation TS PQ CFC ELMO, IELE, PELE}
	\begin{questions}


\question
Pour chaque affirmation ci-dessous concernant le système de bus KNX, évaluer si elle est vraie ou fausse en cochant la case correspondante~?
\begin{parts} 
	\part[1] La tension d'alimentation nominale du bus KNX (TP) fournie par l'alimentation est de \qty{29}{[\volt]} DC. 
	\begin{checkboxes} 
		\correctchoice Vrai 
		\choice Faux 
	\end{checkboxes}
	
	\part[1] Dans une installation KNX TP (Torsadé), la topologie en boucle (anneau) est autorisée pour garantir la redondance.
	\begin{checkboxes}
		\choice Vrai
		\correctchoice Faux
	\end{checkboxes}
	
	\part[1] Une adresse physique se terminant par zéro (exemple : 1.2.0) désigne obligatoirement un coupleur (de ligne ou de zone).
	\begin{checkboxes}
		\correctchoice Vrai
		\choice Faux
	\end{checkboxes}
	
	\part[1] La distance maximale de câble autorisée entre deux participants quelconques raccordés sur une même ligne est de \qty{700}{[\meter]}.
	\begin{checkboxes}
		\correctchoice Vrai
		\choice Faux
	\end{checkboxes}
\end{parts}
\vspace{0.5em}
\needspace{40\baselineskip}
\begin{solution} 
	\begin{align*} 
		\text{a) Tension nominale } U_{bus} &= \qty{29}{[\volt]} \text{ (Tolérance entre } \qty{21}{[\volt]} \text{ et } \qty{30}{[\volt]} \text{)} \implies \text{Vrai} \\[6pt]
		 \text{b) Topologie Anneau} &= \text{Strictement interdite (Risque de collision de télégrammes)} \implies \text{Faux} \\[6pt]
		 \text{c) Adresse physique } X.Y.0 &= \text{Adresse réservée aux coupleurs de ligne ou de zone} \implies \text{Vrai} \\[6pt] 
		 \text{d) Distance max (Appareil-Appareil)} &= \qty{700}{[\meter]} \text{ (et } \qty{350}{[\meter]} \text{ entre l'alim et un appareil)} \implies \text{Vrai}
	\end{align*}
	\begin{verbatim} 
		% Télécharger OCTAVE depuis https://wiki.octave.org/Using_Octave et l'installer 
		% Puis copier-coller le code dans OCTAVE
		
		% Parametres de conception KNX 
		U_bus_nom = 29;              % Tension nominale [V] L_max_alim_dev = 350;
		% Distance max alim-appareil [m] L_max_dev_dev = 700;         
		
		% Distance max appareil-appareil [m] L_tot_ligne = 1000;          
		% Longueur totale de cable par ligne [m]
		% Verification des affirmations printf("--- Analyse de la technologie KNX ---\n"); 
		printf("1. Tension nominale du bus : %.0f [V] DC\n", U_bus_nom); 
		printf("2. Topologie autorisee : Bus, Etoile, Arbre (Boucle interdite).\n"); 
		printf("3. Adresse physique X.Y.0 : Coupleur de ligne (Z.L.0)\n"); 
		printf("4. Distance maximale entre 2 participants : %.0f [m]\n", L_max_dev_dev); 
		printf("5. Distance maximale Alimentation-Participant : %.0f [m]\n", L_max_alim_dev); 
	\end{verbatim}
	\end{solution}
	
\end{questions}
\end{document}
